\documentclass[11pt,letterpaper]{article}

\usepackage[spanish,es-noshorthands,es-tabla]{babel}
\usepackage{fontspec}
% Nimbus Sans es una Helvetica libre: calza con el look & feel del MIT.
\setmainfont{Nimbus Sans}
% Símbolo del colón, que Nimbus Sans no trae.
\newfontfamily\fuentecolon{Liberation Sans}
\newcommand{\colones}{{\fuentecolon ₡}}
\usepackage{microtype}
\usepackage{geometry}
\usepackage{enumitem}
\usepackage{xcolor}
\usepackage{csquotes}
\usepackage{tabularx}
\usepackage{array}
\usepackage{booktabs}
\usepackage{colortbl}
\usepackage{titlesec}
\usepackage{fancyhdr}
\usepackage{needspace}
\usepackage[colorlinks=true,urlcolor=mitRojo,linkcolor=black]{hyperref}
\usepackage[breakable]{tcolorbox}

\geometry{hmargin=2.3cm,top=2.2cm,bottom=2.2cm,headheight=14pt}
\hbadness=10000
\setlength{\emergencystretch}{1em}
\setlength{\parindent}{0pt}
\setlength{\parskip}{0.45em}
\setlist{nosep,topsep=0.3em}
\renewcommand{\arraystretch}{1.3}

% Colores del MIT y uno por cada pieza del prompt.
\definecolor{mitRojo}{HTML}{A31F34}
\definecolor{mitGris}{HTML}{8A8B8C}
\definecolor{fondo}{HTML}{F3F3F3}
\definecolor{pRol}{HTML}{A31F34}
\definecolor{pContexto}{HTML}{1F5C99}
\definecolor{pTarea}{HTML}{2E7D32}
\definecolor{pReglas}{HTML}{B35C00}
\definecolor{pEjemplo}{HTML}{6A4C93}
\definecolor{pFormato}{HTML}{5A5B5C}

% Títulos de sección: rojo MIT, con una línea gris debajo.
\titleformat{\section}{\Large\bfseries\color{mitRojo}}{}{0pt}{}
  [{\color{mitGris}\titlerule[0.6pt]}]
\titlespacing*{\section}{0pt}{1.1em}{0.6em}
\titleformat{\subsection}{\large\bfseries}{}{0pt}{}
\titlespacing*{\subsection}{0pt}{0.8em}{0.3em}

% Encabezado y pie discretos.
\pagestyle{fancy}
\fancyhf{}
\renewcommand{\headrulewidth}{0pt}
\fancyhead[L]{\small\color{mitGris}Teach-back. Ingeniería de prompts}
\fancyhead[R]{\small\color{mitGris}Módulo 4. Actividad 4.1}
\fancyfoot[C]{\small\color{mitGris}\thepage}

% Etiqueta de color con el nombre de una pieza del prompt.
%   \etiqueta{color}{texto}
\newcommand{\etiqueta}[2]{%
  \tcbox[on line,colback=#1,colframe=#1,boxrule=0pt,arc=1.5pt,
    left=3pt,right=3pt,top=1pt,bottom=1pt]{\color{white}\footnotesize\bfseries #2}}

% Una pieza del prompt: etiqueta al margen izquierdo y el texto a la derecha.
%   \pieza{color}{nombre}{texto}
\newcommand{\pieza}[3]{%
  \begin{tcolorbox}[colback=#1!6,colframe=#1,boxrule=0pt,
    leftrule=3pt,arc=0pt,left=6pt,right=6pt,top=3pt,bottom=3pt,
    before skip=3pt,after skip=3pt]
  \begin{minipage}[t]{2.3cm}\raggedright\footnotesize\bfseries
    \textcolor{#1}{\MakeUppercase{#2}}\end{minipage}\hfill
  \begin{minipage}[t]{\dimexpr\linewidth-2.3cm-8pt\relax}
    \small\raggedright #3\end{minipage}
  \end{tcolorbox}}

% Recuadro de un prompt (texto que se le da a la IA).
\newtcolorbox{prompt}[1]{breakable,colback=white,colframe=mitGris,
  boxrule=0.6pt,arc=1.5pt,left=6pt,right=6pt,top=6pt,bottom=6pt,
  title={#1},coltitle=black,colbacktitle=fondo,fonttitle=\small\bfseries}

% Recuadro del mensaje que saldría al cliente.
\newtcolorbox{mensaje}[1]{colback=fondo,colframe=fondo,
  boxrule=0pt,arc=4pt,left=10pt,right=10pt,top=6pt,bottom=6pt,
  title={#1},coltitle=mitGris,colbacktitle=fondo,fonttitle=\footnotesize\bfseries,
  toptitle=2pt,bottomtitle=-2pt}

% Aviso de que el ejemplo es ilustrativo.
\newtcolorbox{aviso}{colback=mitRojo!5,colframe=mitRojo,boxrule=0pt,
  leftrule=3pt,arc=0pt,left=8pt,right=8pt,top=5pt,bottom=5pt}

\begin{document}

% === Portada =================================================================
\begin{titlepage}
\newgeometry{margin=2.5cm}
\raggedright
{\color{mitRojo}\rule{\textwidth}{3pt}\par}
\vspace{0.4cm}
{\small\bfseries\color{mitGris}MIT Professional Education,
Chief Technology Officer Program\par}
\vspace{4.0cm}
{\large Inteligencia Artificial Generativa\\
Aplicada a la Transformación Digital\par}
\vspace{0.4cm}
{\normalsize\color{mitGris}Módulo 4. Actividad 4.1. Teach-back\par}
\vspace{1.4cm}
{\fontsize{30}{36}\selectfont\bfseries\color{mitRojo}
Ingeniería de prompts:\\ qué es y por qué importa\par}
\vspace{0.6cm}
{\large Una explicación para el equipo de\\
customer success managers de Subaxa\par}
\vfill
{\color{mitGris}\rule{\textwidth}{0.6pt}\par}
\vspace{0.3cm}
{\large\bfseries Carlos Luis Rojas Aragonés\par}
\vspace{0.1cm}
{\normalsize\color{mitGris}CTO de Subaxa\quad\textbar\quad 23 de noviembre de 2026\par}
\restoregeometry
\end{titlepage}

% === 1. Nota de contextualización (100 a 150 palabras) =======================
\needspace{8\baselineskip}
\section*{1. Nota de contextualización}

Elegí explicar la ingeniería de prompts: qué es y por qué importa. Un prompt
son las instrucciones por escrito que le damos a la IA para que haga un
trabajo. En Subaxa esto pesa todos los días. Nuestro agente de user
engagement escribe mensajes personalizados por correo y WhatsApp para cada
inversionista, y la calidad de esos mensajes depende del prompt.

El público es el equipo de customer success managers de Subaxa. No son
técnicos, pero aprueban cada mensaje antes de que salga. Son la última
revisión entre la IA y el cliente.

La idea clave que obtuve al prepararlo es que un mensaje genérico casi nunca
es culpa de la IA: es un prompt que dejó decisiones sin tomar. Y las
personas que mejor pueden tomar esas decisiones son los CSM, no yo.

% === 2. Explicación del concepto ==============================================
\needspace{8\baselineskip}
\section*{2. Qué es un prompt y cómo funciona}

\needspace{5\baselineskip}
\subsection*{Qué es}

Un \textbf{prompt} son las instrucciones por escrito que le damos a la IA
para que haga un trabajo. La \textbf{IA generativa} (o \emph{modelo}) es el
programa que escribe el texto a partir de esas instrucciones.

La \textbf{ingeniería de prompts} es el oficio de escribir esas
instrucciones de forma que la IA entregue lo que necesitamos, y no lo que
ella supone.

\needspace{5\baselineskip}
\subsection*{Una forma sencilla de verlo: el brief a un compañero nuevo}

Imaginen que llega alguien nuevo al equipo y le piden: \enquote{escribile a
un cliente sobre un remate}. Va a escribir algo correcto, pero genérico.
No conoce al cliente ni nuestras reglas.

Ahora imaginen que le dicen quién es el cliente, qué busca y qué remate
salió. Le dicen qué no puede prometer. Le muestran un mensaje que ya
funcionó y le piden que sea corto, para WhatsApp. El mensaje sale bien
desde la primera vez.

La IA es ese compañero nuevo: escribe rápido y bien, pero solo sabe lo que
le decimos. \textbf{El prompt es el brief.} Si el brief es vago, el
resultado es vago.

\needspace{5\baselineskip}
\subsection*{Cómo funciona: las seis piezas de un buen prompt}

Un buen prompt tiene seis piezas. Cada una responde una pregunta que el
compañero nuevo nos haría. Junto con la tarea, le pasamos a la IA un
\emph{contexto}: la información que necesita, en este caso lo que sabemos
del usuario y del remate.

\begin{center}
\small
\begin{tabularx}{\textwidth}{@{}>{\raggedright\arraybackslash}p{2.9cm}
  >{\raggedright\arraybackslash}X >{\raggedright\arraybackslash}X@{}}
\toprule
\textbf{Pieza} & \textbf{Qué es, en palabras simples} &
  \textbf{En el agente de user engagement}\\
\midrule
\etiqueta{pRol}{Rol} & Quién \enquote{es} la IA en esta tarea &
  Alguien del equipo de Subaxa que acompaña a inversionistas\\
\etiqueta{pContexto}{Contexto} & Lo que sabemos de la persona &
  Zonas, tipo de propiedad, presupuesto, qué remates vio o guardó, etapa
  del cliente\\
\etiqueta{pTarea}{Tarea} & Qué tiene que hacer, una sola cosa &
  Avisar de un remate nuevo que calza con su perfil\\
\etiqueta{pReglas}{Reglas} & Lo que puede y no puede hacer &
  No prometer ganancias, no inventar datos del remate, invitar a hablar con
  el equipo\\
\etiqueta{pEjemplo}{Ejemplo} & Un mensaje modelo de cómo suena
  \enquote{bien} & Un mensaje corto, cercano y profesional ya aprobado\\
\etiqueta{pFormato}{Formato} & Cómo se entrega &
  Canal (WhatsApp o correo), largo máximo, un solo llamado a la acción\\
\bottomrule
\end{tabularx}
\end{center}

La pieza de reglas incluye \enquote{no inventar datos}. Lo pongo porque la
IA a veces inventa un dato que suena creíble. A eso se le llama
\emph{alucinación}.

\needspace{5\baselineskip}
\subsection*{Por qué importa}

\begin{itemize}
\item \textbf{Ustedes aprueban cada mensaje antes de que salga.} Son la
  última revisión entre la IA y el inversionista.
\item \textbf{Un mensaje bueno o malo casi siempre se explica por el
  prompt.} Si falta el contexto, sale genérico. Si falta una regla, puede
  prometer algo que no debe. Si falta el formato, sale largo o en el tono
  equivocado.
\item \textbf{Es el mismo criterio que ya tienen.} Un CSM sabe cómo se le
  habla a un cliente nuevo y cómo a uno inactivo. El prompt es el lugar
  donde ese criterio queda por escrito para la IA.
\end{itemize}

% === 3. Ejemplo en Subaxa =====================================================
\needspace{8\baselineskip}
\section*{3. Un ejemplo en Subaxa: el agente de user engagement}

\needspace{5\baselineskip}
\subsection*{Qué hace el agente hoy}

El agente está en producción. Escribe mensajes por \textbf{correo} y
\textbf{WhatsApp} para cada usuario de la plataforma. Para personalizarlos
usa cuatro tipos de datos:

\begin{itemize}
\item las preferencias del usuario: zonas, tipo de propiedad y presupuesto;
\item lo que hace en la plataforma: remates que vio, guardó o en los que
  participó;
\item los remates nuevos del Boletín Judicial que calzan con su perfil;
\item la etapa del cliente: nuevo, activo, inactivo, etc.
\end{itemize}

Escribe en dos momentos: cuando pasa algo (aparece un remate que calza, o el
usuario hace o deja de hacer algo) y en un resumen periódico personalizado.
Antes de que un mensaje salga, \textbf{un CSM lo aprueba}.

\begin{aviso}
\small\textbf{Importante:} el ejemplo que sigue es \textbf{ilustrativo y
simplificado}. No es el prompt real del agente. La usuaria, el remate y los
montos son ficticios.
\end{aviso}

\needspace{5\baselineskip}
\subsection*{Un prompt pobre}

\begin{prompt}{Prompt pobre}
\small Escriba un mensaje de WhatsApp para un usuario de Subaxa sobre un
remate nuevo.
\end{prompt}

\begin{mensaje}{Mensaje que produciría (genérico)}
\small ¡Hola! En Subaxa tenemos nuevas oportunidades de inversión en remates
judiciales. ¡No se las pierda! Entre a la plataforma y descubra las mejores
propiedades al mejor precio. ¡Invierta hoy!
\end{mensaje}

\textbf{Qué está mal.} No dice cuál remate. No sabe a quién le habla. Suena
a publicidad masiva. Y \enquote{al mejor precio} es una promesa que no
podemos hacer.

% Las seis piezas van juntas en una misma página.
\needspace{24\baselineskip}
\subsection*{El mismo pedido, con las seis piezas}

Cada recuadro de color es una pieza. La etiqueta al margen dice cuál.

\pieza{pRol}{Rol}{Usted es parte del equipo de Subaxa que acompaña a
inversionistas en remates judiciales en Costa Rica. Escribe de forma
cercana y profesional, tratando de \enquote{usted}.}

\pieza{pContexto}{Contexto del usuario}{%
Nombre: Ana (usuaria ficticia).\\
Etapa: activa. Se registró hace unos meses y todavía no ha participado en
un remate.\\
Busca: casas en Heredia. Presupuesto: hasta \colones 80\,000\,000.\\
Comportamiento: la semana pasada guardó dos casas en Heredia en remate.\\[0.3em]
\textbf{Remate nuevo que calza (ficticio):} casa de 3 habitaciones en
Heredia. Base del remate: \colones 65\,000\,000. Primer remate: 15 de enero
(fecha ficticia).}

\pieza{pTarea}{Tarea}{Avísele a Ana de este remate nuevo y por qué podría
interesarle.}

\pieza{pReglas}{Reglas}{%
Use solo los datos de arriba. No invente nada del remate.\\
No prometa ganancias ni diga que es \enquote{una ganga}.\\
No ponga direcciones exactas.\\
Invítela a ver el detalle o a hablar con el equipo legal y de inversión.}

\pieza{pEjemplo}{Ejemplo de tono}{\enquote{Hola, Luis. Vimos que le interesan
los lotes en Cartago y apareció uno nuevo en el Boletín Judicial. Si gusta,
lo revisamos juntos.} (ficticio)}

\pieza{pFormato}{Formato}{WhatsApp. Máximo 4 líneas. Un solo llamado a la
acción. Sin emojis.}

\begin{mensaje}{Mensaje que produciría (personalizado)}
\small Hola, Ana. Salió en el Boletín Judicial una casa de 3 habitaciones en
Heredia, parecida a las que guardó la semana pasada. La base es de
\colones 65\,000\,000 y el primer remate es el 15 de enero. Si gusta, el
equipo legal y de inversión le acompaña a revisarla antes de decidir.
\end{mensaje}

\textbf{Qué cambió.} El mensaje usa su zona, su tipo de propiedad y lo que
hizo en la plataforma. Está dentro de su presupuesto. No promete nada. Y es
corto. La IA es la misma en los dos casos: lo único que cambió fue el
prompt.

% === 4. Conexión con resultados reales =======================================
\needspace{8\baselineskip}
\section*{4. Cómo esto cambia lo que sale y lo que decidimos}

\textbf{Lo que le llega al cliente.} Cada mensaje que el agente manda por
correo o WhatsApp sale de un prompt. Si al prompt le falta contexto, el
inversionista recibe algo genérico, que se parece a cualquier publicidad. Si
le falta una regla, el mensaje puede prometer algo que Subaxa no puede
cumplir. Si le falta el formato, llega un texto largo a un canal donde se
lee rápido. Con las seis piezas, el mensaje habla de lo que esa persona
busca, con los datos correctos del remate.

\textbf{El trabajo de aprobación de los CSM.} Cuando el prompt es pobre,
el CSM termina reescribiendo mensajes en lugar de revisarlos. Cuando el
prompt es bueno, aprobar es más rápido y la revisión se enfoca en lo que de
verdad requiere criterio humano: el tono con un cliente en particular, o un
caso que el prompt no previó.

\textbf{La decisión de fondo: corregir cada mensaje o ajustar el prompt
una vez.} Corregir a mano arregla un mensaje. El mismo error vuelve a
aparecer en el siguiente. Ajustar el prompt arregla la causa, y ese ajuste
mejora todos los mensajes que vienen. Por eso, cuando un CSM ve que un
error se repite, lo más útil no es solo corregirlo: es decir qué pieza
falló y proponer el cambio.

En resumen, los cuatro mensajes que quiero que se lleven:

\begin{enumerate}
\item Un prompt es el brief que le damos a la IA: si el brief es vago, el
  resultado es vago.
\item Un buen prompt tiene seis piezas: rol, contexto del usuario, tarea,
  reglas, ejemplo y formato.
\item Cuando un mensaje sale mal, casi siempre falta o falla una pieza.
\item El CSM no solo aprueba: puede mejorar el prompt, y eso mejora todos
  los mensajes que vienen.
\end{enumerate}

% === 5. Conclusión práctica ===================================================
\needspace{8\baselineskip}
\section*{5. Qué pueden hacer desde mañana}

Al aprobar un mensaje, háganse estas cinco preguntas rápidas:

\begin{tcolorbox}[colback=white,colframe=mitRojo,boxrule=0.8pt,
  arc=1.5pt,left=8pt,right=8pt,top=6pt,bottom=6pt]
\begin{enumerate}[label=\textbf{\color{mitRojo}\arabic*.},leftmargin=1.6em,itemsep=0.3em]
\item ¿Se nota que es para esta persona? (zona, tipo de propiedad, lo que
  hizo)
\item ¿Los datos del remate son correctos y no hay nada inventado?
\item ¿Promete algo que no debe? (ganancias, \enquote{ganga}, certeza del
  resultado)
\item ¿El tono calza con la etapa del cliente? (nuevo, activo, inactivo)
\item ¿Tiene el largo del canal y un solo llamado a la acción?
\end{enumerate}
\end{tcolorbox}

Si una respuesta es \enquote{no}, además de corregir el mensaje, escriban
una sugerencia de ajuste con esta forma simple:

\begin{tcolorbox}[colback=fondo,colframe=fondo,arc=1.5pt,
  left=8pt,right=8pt,top=6pt,bottom=6pt,
  title={Plantilla para sugerir un ajuste al prompt},coltitle=black,
  colbacktitle=fondo,fonttitle=\small\bfseries]
\small
\begin{tabularx}{\linewidth}{@{}>{\bfseries}p{3.6cm}X@{}}
Pieza: & \etiqueta{pReglas}{Reglas}\\
Qué pasó: & El mensaje dijo \enquote{excelente oportunidad de ganancia}.\\
Ajuste propuesto: & Agregar la regla \enquote{no hablar de ganancias ni de
  rentabilidad esperada}.\\
\end{tabularx}
\end{tcolorbox}

Ustedes ya tienen el criterio. Esta plantilla es la forma de dejarlo por
escrito para la IA.

\end{document}
